\begin{register}{H}{UART\_FIFO\_REG}{0x0}\label{regdesc:UARTFIFOREG}
\regfield{(reserved)}{24}{8}{000000000000000000000000}%
\regfield{UART\_RXFIFO\_RD\_BYTE}{8}{0}{00000000}%
\reglabel{Reset}\regnewline%
\begin{regdesc}\begin{reglist}
\item [UART\_RXFIFO\_RD\_BYTE] UART\textcolor{red}{\it{n}} accesses FIFO via this register. (R/W)
\end{reglist}\end{regdesc}
\end{register}


\begin{register}{H}{UART\_INT\_RAW\_REG}{0x4}\label{regdesc:UARTINTRAWREG}
\regfield{(reserved)}{13}{19}{0000000000000}%
\regfield{UART\_AT\_CMD\_CHAR\_DET\_INT\_RAW}{1}{18}{0}%
\regfield{UART\_RS485\_CLASH\_INT\_RAW}{1}{17}{0}%
\regfield{UART\_RS485\_FRM\_ERR\_INT\_RAW}{1}{16}{0}%
\regfield{UART\_RS485\_PARITY\_ERR\_INT\_RAW}{1}{15}{0}%
\regfield{UART\_TX\_DONE\_INT\_RAW}{1}{14}{0}%
\regfield{UART\_TX\_BRK\_IDLE\_DONE\_INT\_RAW}{1}{13}{0}%
\regfield{UART\_TX\_BRK\_DONE\_INT\_RAW}{1}{12}{0}%
\regfield{UART\_GLITCH\_DET\_INT\_RAW}{1}{11}{0}%
\regfield{UART\_SW\_XOFF\_INT\_RAW}{1}{10}{0}%
\regfield{UART\_SW\_XON\_INT\_RAW}{1}{9}{0}%
\regfield{UART\_RXFIFO\_TOUT\_INT\_RAW}{1}{8}{0}%
\regfield{UART\_BRK\_DET\_INT\_RAW}{1}{7}{0}%
\regfield{UART\_CTS\_CHG\_INT\_RAW}{1}{6}{0}%
\regfield{UART\_DSR\_CHG\_INT\_RAW}{1}{5}{0}%
\regfield{UART\_RXFIFO\_OVF\_INT\_RAW}{1}{4}{0}%
\regfield{UART\_FRM\_ERR\_INT\_RAW}{1}{3}{0}%
\regfield{UART\_PARITY\_ERR\_INT\_RAW}{1}{2}{0}%
\regfield{UART\_TXFIFO\_EMPTY\_INT\_RAW}{1}{1}{0}%
\regfield{UART\_RXFIFO\_FULL\_INT\_RAW}{1}{0}{0}%
\reglabel{Reset}\regnewline%
\begin{regdesc}\begin{reglist}
\item [UART\_AT\_CMD\_CHAR\_DET\_INT\_RAW]The raw interrupt status bit for the \hyperref[int:UARTATCMDCHARDETINT]{UART\_AT\_CMD\_CHAR\_DET\_INT} interrupt. (RO)
\item [UART\_RS485\_CLASH\_INT\_RAW]The raw interrupt status bit for the \hyperref[int:UARTRS485CLASHINT]{UART\_RS485\_CLASH\_INT} interrupt. (RO)
\item [UART\_RS485\_FRM\_ERR\_INT\_RAW]The raw interrupt status bit for the \hyperref[int:UARTRS485FRMERRINT]{UART\_RS485\_FRM\_ERR\_INT} interrupt. (RO)
\item [UART\_RS485\_PARITY\_ERR\_INT\_RAW]The raw interrupt status bit for the \hyperref[int:UARTRS485PARITYERRINT]{UART\_RS485\_PARITY\_ERR\_INT} interrupt. (RO)
\item [UART\_TX\_DONE\_INT\_RAW]The raw interrupt status bit for the \hyperref[int:UARTTXDONEINT]{UART\_TX\_DONE\_INT} interrupt. (RO)
\item [UART\_TX\_BRK\_IDLE\_DONE\_INT\_RAW]The raw interrupt status bit for the \hyperref[int:UARTTXBRKIDLEDONEINT]{UART\_TX\_BRK\_IDLE\_DONE\_INT} interrupt. (RO)
\item [UART\_TX\_BRK\_DONE\_INT\_RAW]The raw interrupt status bit for the \hyperref[int:UARTTXBRKDONEINT]{UART\_TX\_BRK\_DONE\_INT} interrupt. (RO)
\item [UART\_GLITCH\_DET\_INT\_RAW]The raw interrupt status bit for the \hyperref[int:UARTGLITCHDETINT]{UART\_GLITCH\_DET\_INT} interrupt. (RO)
\item [UART\_SW\_XOFF\_INT\_RAW]The raw interrupt status bit for the \hyperref[int:UARTSWXOFFINT]{UART\_SW\_XOFF\_INT} interrupt. (RO)
\item [UART\_SW\_XON\_INT\_RAW]The raw interrupt status bit for the \hyperref[int:UARTSWXONINT]{UART\_SW\_XON\_INT} interrupt. (RO)
\item [UART\_RXFIFO\_TOUT\_INT\_RAW]The raw interrupt status bit for the \hyperref[int:UARTRXFIFOTOUTINT]{UART\_RXFIFO\_TOUT\_INT} interrupt. (RO)
\item [UART\_BRK\_DET\_INT\_RAW]The raw interrupt status bit for the \hyperref[int:UARTBRKDETINT]{UART\_BRK\_DET\_INT} interrupt. (RO)
\item [UART\_CTS\_CHG\_INT\_RAW]The raw interrupt status bit for the \hyperref[int:UARTCTSCHGINT]{UART\_CTS\_CHG\_INT} interrupt. (RO)
\item [UART\_DSR\_CHG\_INT\_RAW]The raw interrupt status bit for the \hyperref[int:UARTDSRCHGINT]{UART\_DSR\_CHG\_INT} interrupt. (RO)
\item [UART\_RXFIFO\_OVF\_INT\_RAW]The raw interrupt status bit for the \hyperref[int:UARTRXFIFOOVFINT]{UART\_RXFIFO\_OVF\_INT} interrupt. (RO)
\item [UART\_FRM\_ERR\_INT\_RAW]The raw interrupt status bit for the \hyperref[int:UARTFRMERRINT]{UART\_FRM\_ERR\_INT} interrupt. (RO)
\item [Continued on the next page...]
\end{reglist}\end{regdesc}
\end{register}
\addtocounter{Regfloat}{-1}
\begin{register}{H}{UART\_INT\_RAW\_REG}{0x4}
\begin{regdesc}\begin{reglist}
\item [Continued from the previous page...]
\item [UART\_PARITY\_ERR\_INT\_RAW]The raw interrupt status bit for the \hyperref[int:UARTPARITYERRINT]{UART\_PARITY\_ERR\_INT} interrupt. (RO)
\item [UART\_TXFIFO\_EMPTY\_INT\_RAW]The raw interrupt status bit for the \hyperref[int:UARTTXFIFOEMPTYINT]{UART\_TXFIFO\_EMPTY\_INT} interrupt. (RO)
\item [UART\_RXFIFO\_FULL\_INT\_RAW]The raw interrupt status bit for the \hyperref[int:UARTRXFIFOFULLINT]{UART\_RXFIFO\_FULL\_INT} interrupt. (RO)
\end{reglist}\end{regdesc}
\end{register}


\begin{register}{H}{UART\_INT\_ST\_REG}{0x8}\label{regdesc:UARTINTSTREG}
\regfield{(reserved)}{13}{19}{0000000000000}%
\regfield{UART\_AT\_CMD\_CHAR\_DET\_INT\_ST}{1}{18}{0}%
\regfield{UART\_RS485\_CLASH\_INT\_ST}{1}{17}{0}%
\regfield{UART\_RS485\_FRM\_ERR\_INT\_ST}{1}{16}{0}%
\regfield{UART\_RS485\_PARITY\_ERR\_INT\_ST}{1}{15}{0}%
\regfield{UART\_TX\_DONE\_INT\_ST}{1}{14}{0}%
\regfield{UART\_TX\_BRK\_IDLE\_DONE\_INT\_ST}{1}{13}{0}%
\regfield{UART\_TX\_BRK\_DONE\_INT\_ST}{1}{12}{0}%
\regfield{UART\_GLITCH\_DET\_INT\_ST}{1}{11}{0}%
\regfield{UART\_SW\_XOFF\_INT\_ST}{1}{10}{0}%
\regfield{UART\_SW\_XON\_INT\_ST}{1}{9}{0}%
\regfield{UART\_RXFIFO\_TOUT\_INT\_ST}{1}{8}{0}%
\regfield{UART\_BRK\_DET\_INT\_ST}{1}{7}{0}%
\regfield{UART\_CTS\_CHG\_INT\_ST}{1}{6}{0}%
\regfield{UART\_DSR\_CHG\_INT\_ST}{1}{5}{0}%
\regfield{UART\_RXFIFO\_OVF\_INT\_ST}{1}{4}{0}%
\regfield{UART\_FRM\_ERR\_INT\_ST}{1}{3}{0}%
\regfield{UART\_PARITY\_ERR\_INT\_ST}{1}{2}{0}%
\regfield{UART\_TXFIFO\_EMPTY\_INT\_ST}{1}{1}{0}%
\regfield{UART\_RXFIFO\_FULL\_INT\_ST}{1}{0}{0}%
\reglabel{Reset}\regnewline%
\begin{regdesc}\begin{reglist}
\item [UART\_AT\_CMD\_CHAR\_DET\_INT\_ST]The masked interrupt status bit for the \hyperref[int:UARTATCMDCHARDETINT]{UART\_AT\_CMD\_CHAR\_DET\_INT} interrupt. (RO)
\item [UART\_RS485\_CLASH\_INT\_ST]The masked interrupt status bit for the \hyperref[int:UARTRS485CLASHINT]{UART\_RS485\_CLASH\_INT} interrupt. (RO)
\item [UART\_RS485\_FRM\_ERR\_INT\_ST]The masked interrupt status bit for the \hyperref[int:UARTRS485FRMERRINT]{UART\_RS485\_FRM\_ERR\_INT} interrupt. (RO)
\item [UART\_RS485\_PARITY\_ERR\_INT\_ST]The masked interrupt status bit for the \hyperref[int:UARTRS485PARITYERRINT]{UART\_RS485\_PARITY\_ERR\_INT} interrupt. (RO)
\item [UART\_TX\_DONE\_INT\_ST]The masked interrupt status bit for the \hyperref[int:UARTTXDONEINT]{UART\_TX\_DONE\_INT} interrupt. (RO)
\item [UART\_TX\_BRK\_IDLE\_DONE\_INT\_ST]The masked interrupt status bit for the \hyperref[int:UARTTXBRKIDLEDONEINT]{UART\_TX\_BRK\_IDLE\_DONE\_INT} interrupt. (RO)
\item [UART\_TX\_BRK\_DONE\_INT\_ST]The masked interrupt status bit for the \hyperref[int:UARTTXBRKDONEINT]{UART\_TX\_BRK\_DONE\_INT} interrupt. (RO)
\item [UART\_GLITCH\_DET\_INT\_ST]The masked interrupt status bit for the \hyperref[int:UARTGLITCHDETINT]{UART\_GLITCH\_DET\_INT} interrupt. (RO)
\item [UART\_SW\_XOFF\_INT\_ST]The masked interrupt status bit for the \hyperref[int:UARTSWXOFFINT]{UART\_SW\_XOFF\_INT} interrupt. (RO)
\item [UART\_SW\_XON\_INT\_ST]The masked interrupt status bit for the \hyperref[int:UARTSWXONINT]{UART\_SW\_XON\_INT} interrupt. (RO)
\item [UART\_RXFIFO\_TOUT\_INT\_ST]The masked interrupt status bit for the \hyperref[int:UARTRXFIFOTOUTINT]{UART\_RXFIFO\_TOUT\_INT} interrupt. (RO)
\item [UART\_BRK\_DET\_INT\_ST]The masked interrupt status bit for the \hyperref[int:UARTBRKDETINT]{UART\_BRK\_DET\_INT} interrupt. (RO)
\item [UART\_CTS\_CHG\_INT\_ST]The masked interrupt status bit for the \hyperref[int:UARTCTSCHGINT]{UART\_CTS\_CHG\_INT} interrupt. (RO)
\item [UART\_DSR\_CHG\_INT\_ST]The masked interrupt status bit for the \hyperref[int:UARTDSRCHGINT]{UART\_DSR\_CHG\_INT} interrupt. (RO)
\item [Continued on the next page...]
\end{reglist}\end{regdesc}
\end{register}

\addtocounter{Regfloat}{-1}
\begin{register}{H}{UART\_INT\_ST\_REG}{0x8}
\begin{regdesc}\begin{reglist}
\item [Continued from the previous page...]
\item [UART\_RXFIFO\_OVF\_INT\_ST]The masked interrupt status bit for the \hyperref[int:UARTRXFIFOOVFINT]{UART\_RXFIFO\_OVF\_INT} interrupt. (RO)
\item [UART\_FRM\_ERR\_INT\_ST]The masked interrupt status bit for the \hyperref[int:UARTFRMERRINT]{UART\_FRM\_ERR\_INT} interrupt. (RO)
\item [UART\_PARITY\_ERR\_INT\_ST]The masked interrupt status bit for the \hyperref[int:UARTPARITYERRINT]{UART\_PARITY\_ERR\_INT} interrupt. (RO)
\item [UART\_TXFIFO\_EMPTY\_INT\_ST]The masked interrupt status bit for the \hyperref[int:UARTTXFIFOEMPTYINT]{UART\_TXFIFO\_EMPTY\_INT} interrupt. (RO)
\item [UART\_RXFIFO\_FULL\_INT\_ST]The masked interrupt status bit for \hyperref[int:UARTRXFIFOFULLINT]{UART\_RXFIFO\_FULL\_INT}. (RO)
\end{reglist}\end{regdesc}
\end{register}


\begin{register}{H}{UART\_INT\_ENA\_REG}{0xC}\label{regdesc:UARTINTENAREG}
\regfield{(reserved)}{13}{19}{0000000000000}%
\regfield{UART\_AT\_CMD\_CHAR\_DET\_INT\_ENA}{1}{18}{0}%
\regfield{UART\_RS485\_CLASH\_INT\_ENA}{1}{17}{0}%
\regfield{UART\_RS485\_FRM\_ERR\_INT\_ENA}{1}{16}{0}%
\regfield{UART\_RS485\_PARITY\_ERR\_INT\_ENA}{1}{15}{0}%
\regfield{UART\_TX\_DONE\_INT\_ENA}{1}{14}{0}%
\regfield{UART\_TX\_BRK\_IDLE\_DONE\_INT\_ENA}{1}{13}{0}%
\regfield{UART\_TX\_BRK\_DONE\_INT\_ENA}{1}{12}{0}%
\regfield{UART\_GLITCH\_DET\_INT\_ENA}{1}{11}{0}%
\regfield{UART\_SW\_XOFF\_INT\_ENA}{1}{10}{0}%
\regfield{UART\_SW\_XON\_INT\_ENA}{1}{9}{0}%
\regfield{UART\_RXFIFO\_TOUT\_INT\_ENA}{1}{8}{0}%
\regfield{UART\_BRK\_DET\_INT\_ENA}{1}{7}{0}%
\regfield{UART\_CTS\_CHG\_INT\_ENA}{1}{6}{0}%
\regfield{UART\_DSR\_CHG\_INT\_ENA}{1}{5}{0}%
\regfield{UART\_RXFIFO\_OVF\_INT\_ENA}{1}{4}{0}%
\regfield{UART\_FRM\_ERR\_INT\_ENA}{1}{3}{0}%
\regfield{UART\_PARITY\_ERR\_INT\_ENA}{1}{2}{0}%
\regfield{UART\_TXFIFO\_EMPTY\_INT\_ENA}{1}{1}{0}%
\regfield{UART\_RXFIFO\_FULL\_INT\_ENA}{1}{0}{0}%
\reglabel{Reset}\regnewline%
\begin{regdesc}\begin{reglist}
\item [UART\_AT\_CMD\_CHAR\_DET\_INT\_ENA]The interrupt enable bit for the \hyperref[int:UARTATCMDCHARDETINT]{UART\_AT\_CMD\_CHAR\_DET\_INT} interrupt. (R/W)
\item [UART\_RS485\_CLASH\_INT\_ENA]The interrupt enable bit for the \hyperref[int:UARTRS485CLASHINT]{UART\_RS485\_CLASH\_INT} interrupt. (R/W)
\item [UART\_RS485\_FRM\_ERR\_INT\_ENA]The interrupt enable bit for the \hyperref[int:UARTRS485FRMERRINT]{UART\_RS485\_FRM\_ERR\_INT} interrupt. (R/W)
\item [UART\_RS485\_PARITY\_ERR\_INT\_ENA]The interrupt enable bit for the \hyperref[int:UARTRS485PARITYERRINT]{UART\_RS485\_PARITY\_ERR\_INT} interrupt. (R/W)
\item [UART\_TX\_DONE\_INT\_ENA]The interrupt enable bit for the \hyperref[int:UARTTXDONEINT]{UART\_TX\_DONE\_INT} interrupt. (R/W)
\item [UART\_TX\_BRK\_IDLE\_DONE\_INT\_ENA]The interrupt enable bit for the \hyperref[int:UARTTXBRKIDLEDONEINT]{UART\_TX\_BRK\_IDLE\_DONE\_INT} interrupt. (R/W)
\item [UART\_TX\_BRK\_DONE\_INT\_ENA]The interrupt enable bit for the \hyperref[int:UARTTXBRKDONEINT]{UART\_TX\_BRK\_DONE\_INT} interrupt. (R/W)
\item [UART\_GLITCH\_DET\_INT\_ENA]The interrupt enable bit for the \hyperref[int:UARTGLITCHDETINT]{UART\_GLITCH\_DET\_INT} interrupt. (R/W)
\item [UART\_SW\_XOFF\_INT\_ENA]The interrupt enable bit for the \hyperref[int:UARTSWXOFFINT]{UART\_SW\_XOFF\_INT} interrupt. (R/W)
\item [UART\_SW\_XON\_INT\_ENA]The interrupt enable bit for the \hyperref[int:UARTSWXONINT]{UART\_SW\_XON\_INT} interrupt. (R/W)
\item [UART\_RXFIFO\_TOUT\_INT\_ENA]The interrupt enable bit for the \hyperref[int:UARTRXFIFOTOUTINT]{UART\_RXFIFO\_TOUT\_INT} interrupt. (R/W)
\item [UART\_BRK\_DET\_INT\_ENA]The interrupt enable bit for the \hyperref[int:UARTBRKDETINT]{UART\_BRK\_DET\_INT} interrupt. (R/W)
\item [UART\_CTS\_CHG\_INT\_ENA]The interrupt enable bit for the \hyperref[int:UARTCTSCHGINT]{UART\_CTS\_CHG\_INT} interrupt. (R/W)
\item [UART\_DSR\_CHG\_INT\_ENA]The interrupt enable bit for the \hyperref[int:UARTDSRCHGINT]{UART\_DSR\_CHG\_INT} interrupt. (R/W)
\item [UART\_RXFIFO\_OVF\_INT\_ENA]The interrupt enable bit for the \hyperref[int:UARTRXFIFOOVFINT]{UART\_RXFIFO\_OVF\_INT} interrupt. (R/W)
\item [UART\_FRM\_ERR\_INT\_ENA]The interrupt enable bit for the \hyperref[int:UARTFRMERRINT]{UART\_FRM\_ERR\_INT} interrupt. (R/W)
\item [UART\_PARITY\_ERR\_INT\_ENA]The interrupt enable bit for the \hyperref[int:UARTPARITYERRINT]{UART\_PARITY\_ERR\_INT} interrupt. (R/W)
\item [Continued on the next page...]
\end{reglist}\end{regdesc}
\end{register}
\addtocounter{Regfloat}{-1}
\begin{register}{H}{UART\_INT\_ENA\_REG}{0xC}
\begin{regdesc}\begin{reglist}
\item [Continued from the previous page...]
\item [UART\_TXFIFO\_EMPTY\_INT\_ENA]The interrupt enable bit for the \hyperref[int:UARTTXFIFOEMPTYINT]{UART\_TXFIFO\_EMPTY\_INT} interrupt. (R/W)
\item [UART\_RXFIFO\_FULL\_INT\_ENA]The interrupt enable bit for the \hyperref[int:UARTRXFIFOFULLINT]{UART\_RXFIFO\_FULL\_INT} interrupt. (R/W)
\end{reglist}\end{regdesc}
\end{register}

\begin{register}{H}{UART\_INT\_CLR\_REG}{0x10}\label{regdesc:UARTINTCLRREG}
\regfield{(reserved)}{13}{19}{0000000000000}%
\regfield{UART\_AT\_CMD\_CHAR\_DET\_INT\_CLR}{1}{18}{0}%
\regfield{UART\_RS485\_CLASH\_INT\_CLR}{1}{17}{0}%
\regfield{UART\_RS485\_FRM\_ERR\_INT\_CLR}{1}{16}{0}%
\regfield{UART\_RS485\_PARITY\_ERR\_INT\_CLR}{1}{15}{0}%
\regfield{UART\_TX\_DONE\_INT\_CLR}{1}{14}{0}%
\regfield{UART\_TX\_BRK\_IDLE\_DONE\_INT\_CLR}{1}{13}{0}%
\regfield{UART\_TX\_BRK\_DONE\_INT\_CLR}{1}{12}{0}%
\regfield{UART\_GLITCH\_DET\_INT\_CLR}{1}{11}{0}%
\regfield{UART\_SW\_XOFF\_INT\_CLR}{1}{10}{0}%
\regfield{UART\_SW\_XON\_INT\_CLR}{1}{9}{0}%
\regfield{UART\_RXFIFO\_TOUT\_INT\_CLR}{1}{8}{0}%
\regfield{UART\_BRK\_DET\_INT\_CLR}{1}{7}{0}%
\regfield{UART\_CTS\_CHG\_INT\_CLR}{1}{6}{0}%
\regfield{UART\_DSR\_CHG\_INT\_CLR}{1}{5}{0}%
\regfield{UART\_RXFIFO\_OVF\_INT\_CLR}{1}{4}{0}%
\regfield{UART\_FRM\_ERR\_INT\_CLR}{1}{3}{0}%
\regfield{UART\_PARITY\_ERR\_INT\_CLR}{1}{2}{0}%
\regfield{UART\_TXFIFO\_EMPTY\_INT\_CLR}{1}{1}{0}%
\regfield{UART\_RXFIFO\_FULL\_INT\_CLR}{1}{0}{0}%
\reglabel{Reset}\regnewline%
\begin{regdesc}\begin{reglist}
\item [UART\_AT\_CMD\_CHAR\_DET\_INT\_CLR]Set this bit to clear the \hyperref[int:UARTATCMDCHARDETINT]{UART\_AT\_CMD\_CHAR\_DET\_INT} interrupt. (WO)
\item [UART\_RS485\_CLASH\_INT\_CLR]Set this bit to clear the \hyperref[int:UARTRS485CLASHINT]{UART\_RS485\_CLASH\_INT} interrupt. (WO)
\item [UART\_RS485\_FRM\_ERR\_INT\_CLR]Set this bit to clear the \hyperref[int:UARTRS485FRMERRINT]{UART\_RS485\_FRM\_ERR\_INT} interrupt. (WO)
\item [UART\_RS485\_PARITY\_ERR\_INT\_CLR]Set this bit to clear the \hyperref[int:UARTRS485PARITYERRINT]{UART\_RS485\_PARITY\_ERR\_INT} interrupt. (WO)
\item [UART\_TX\_DONE\_INT\_CLR]Set this bit to clear the \hyperref[int:UARTTXDONEINT]{UART\_TX\_DONE\_INT} interrupt. (WO)
\item [UART\_TX\_BRK\_IDLE\_DONE\_INT\_CLR]Set this bit to clear the \hyperref[int:UARTTXBRKIDLEDONEINT]{UART\_TX\_BRK\_IDLE\_DONE\_INT} interrupt. (WO)
\item [UART\_TX\_BRK\_DONE\_INT\_CLR]Set this bit to clear the \hyperref[int:UARTTXBRKDONEINT]{UART\_TX\_BRK\_DONE\_INT} interrupt. (WO)
\item [UART\_GLITCH\_DET\_INT\_CLR]Set this bit to clear the \hyperref[int:UARTGLITCHDETINT]{UART\_GLITCH\_DET\_INT} interrupt. (WO)
\item [UART\_SW\_XOFF\_INT\_CLR]Set this bit to clear the \hyperref[int:UARTSWXOFFINT]{UART\_SW\_XOFF\_INT} interrupt. (WO)
\item [UART\_SW\_XON\_INT\_CLR]Set this bit to clear the \hyperref[int:UARTSWXONINT]{UART\_SW\_XON\_INT} interrupt. (WO)
\item [UART\_RXFIFO\_TOUT\_INT\_CLR]\label{UARTRXFIFOTOUTINTCLR} Set this bit to clear the \hyperref[int:UARTRXFIFOTOUTINT]{UART\_RXFIFO\_TOUT\_INT} interrupt. This bit can be set only when both rxfifo\_cnt and rx\_mem\_cnt are 0. (WO)
\item [UART\_BRK\_DET\_INT\_CLR]Set this bit to clear the \hyperref[int:UARTBRKDETINT]{UART\_BRK\_DET\_INT} interrupt. (WO)
\item [UART\_CTS\_CHG\_INT\_CLR]Set this bit to clear the \hyperref[int:UARTCTSCHGINT]{UART\_CTS\_CHG\_INT} interrupt. (WO)
\item [UART\_DSR\_CHG\_INT\_CLR]Set this bit to clear the \hyperref[int:UARTDSRCHGINT]{UART\_DSR\_CHG\_INT} interrupt. (WO)
\item [UART\_RXFIFO\_OVF\_INT\_CLR]Set this bit to clear the \hyperref[int:UARTRXFIFOOVFINT]{UART\_RXFIFO\_OVF\_INT} interrupt. (WO)
\item [UART\_FRM\_ERR\_INT\_CLR]Set this bit to clear the \hyperref[int:UARTFRMERRINT]{UART\_FRM\_ERR\_INT} interrupt. (WO)
\item [UART\_PARITY\_ERR\_INT\_CLR]Set this bit to clear the \hyperref[int:UARTPARITYERRINT]{UART\_PARITY\_ERR\_INT} interrupt. (WO)
\item [UART\_TXFIFO\_EMPTY\_INT\_CLR]Set this bit to clear the \hyperref[int:UARTTXFIFOEMPTYINT]{UART\_TXFIFO\_EMPTY\_INT} interrupt. (WO)
\item [UART\_RXFIFO\_FULL\_INT\_CLR]\label{UARTRXFIFOFULLINTCLR} Set this bit to clear the \hyperref[int:UARTRXFIFOFULLINT]{UART\_RXFIFO\_FULL\_INT} interrupt. This bit can be set only when data in Rx\_FIFO is less than \hyperref[UARTRXFIFOFULLTHRHD]{UART\_RXFIFO\_FULL\_THRHD}. (WO)
\end{reglist}\end{regdesc}
\end{register}


\begin{register}{H}{UART\_CLKDIV\_REG}{0x14}\label{regdesc:UARTCLKDIVREG}
\regfield{(reserved)}{8}{24}{00000000}%
\regfield{UART\_CLKDIV\_FRAG}{4}{20}{{0x00}}%
\regfield{UART\_CLKDIV}{20}{0}{{0x0002B6}}%
\reglabel{Reset}\regnewline%
\begin{regdesc}\begin{reglist}
\item [UART\_CLKDIV\_FRAG]The decimal part of the frequency divider factor. (R/W)
\item [UART\_CLKDIV]The integral part of the frequency divider factor. (R/W)
\end{reglist}\end{regdesc}
\end{register}


\begin{register}{H}{UART\_AUTOBAUD\_REG}{0x18}\label{regdesc:UARTAUTOBAUDREG}
\regfield{(reserved)}{16}{16}{0000000000000000}%
\regfield{UART\_GLITCH\_FILT}{8}{8}{{0x010}}%
\regfield{(reserved)}{7}{1}{0000000}%
\regfield{UART\_AUTOBAUD\_EN}{1}{0}{0}%
\reglabel{Reset}\regnewline%
\begin{regdesc}\begin{reglist}
\item [UART\_GLITCH\_FILT]When the input pulse width is lower than this value, the pulse is ignored. This register is used in the autobauding process. (R/W)
\item [UART\_AUTOBAUD\_EN]This is the enable bit for autobaud. (R/W)
\end{reglist}\end{regdesc}
\end{register}


\begin{register}{H}{UART\_STATUS\_REG}{0x1C}\label{regdesc:UARTSTATUSREG}
\regfield{UART\_TXD}{1}{31}{{0}}%
\regfield{UART\_RTSN}{1}{30}{0}%
\regfield{UART\_DTRN}{1}{29}{0}%
\regfield{(reserved)}{1}{28}{0}%
\regfield{UART\_ST\_UTX\_OUT}{4}{24}{0000}%
\regfield{UART\_TXFIFO\_CNT}{8}{16}{00000000}%
\regfield{UART\_RXD}{1}{15}{0}%
\regfield{UART\_CTSN}{1}{14}{0}%
\regfield{UART\_DSRN}{1}{13}{0}%
\regfield{(reserved)}{1}{12}{0}%
\regfield{UART\_ST\_URX\_OUT}{4}{8}{0000}%
\regfield{UART\_RXFIFO\_CNT}{8}{0}{00000000}%
\reglabel{Reset}\regnewline%
\begin{regdesc}\begin{reglist}
\item [UART\_TXD]This bit represents the level of the internal UART RxD signal. (RO)
\item [UART\_RTSN]This bit corresponds to the level of the internal UART CTS signal. (RO)
\item [UART\_DTRN]This bit corresponds to the level of the internal UAR DSR signal. (RO)
\item [UART\_ST\_UTX\_OUT]This register stores the state of the transmitter's finite state machine. 0: TX\_IDLE; 1: TX\_STRT; 2: TX\_DAT0; 3: TX\_DAT1; 4: TX\_DAT2; 5: TX\_DAT3; 6: TX\_DAT4; 7: TX\_DAT5; 8: TX\_DAT6; 9: TX\_DAT7; 10: TX\_PRTY; 11: TX\_STP1; 12: TX\_STP2; 13: TX\_DL0;  14: TX\_DL1. (RO)
\item [UART\_TXFIFO\_CNT]\label{UARTTXFIFOCNT}(tx\_mem\_cnt, txfifo\_cnt) stores the number of bytes of valid data in transmit-FIFO. tx\_mem\_cnt stores the three most significant bits, txfifo\_cnt stores the eight least significant bits. (RO)
\item [UART\_RXD]This bit corresponds to the level of the internal UART RxD signal. (RO)
\item [UART\_CTSN]This bit corresponds to the level of the internal UART CTS signal. (RO)
\item [UART\_DSRN]This bit corresponds to the level of the internal UAR DSR signal. (RO)
\item [UART\_ST\_URX\_OUT]This register stores the value of the receiver's finite state machine. 0: RX\_IDLE; 1: RX\_STRT; 2: RX\_DAT0; 3: RX\_DAT1; 4: RX\_DAT2; 5: RX\_DAT3; 6: RX\_DAT4; 7: RX\_DAT5; 8: RX\_DAT6; 9: RX\_DAT7; 10: RX\_PRTY; 11: RX\_STP1; 12:RX\_STP2; 13: RX\_DL1. (RO)
\item [UART\_RXFIFO\_CNT]\label{UARTRXFIFOCNT}(rx\_mem\_cnt, rxfifo\_cnt) stores the number of bytes of valid data in the receive-FIFO. rx\_mem\_cnt register stores the three most significant bits, rxfifo\_cnt stores the eight least significant bits. (RO)
\end{reglist}\end{regdesc}
\end{register}

\begin{register}{H}{UART\_CONF0\_REG}{0x20}\label{regdesc:UARTCONF0REG}
\regfield{(reserved)}{4}{28}{0000}%
\regfield{UART\_TICK\_REF\_ALWAYS\_ON}{1}{27}{1}%
\regfield{(reserved)}{2}{25}{00}%
\regfield{UART\_DTR\_INV}{1}{24}{0}%
\regfield{UART\_RTS\_INV}{1}{23}{0}%
\regfield{UART\_TXD\_INV}{1}{22}{0}%
\regfield{UART\_DSR\_INV}{1}{21}{0}%
\regfield{UART\_CTS\_INV}{1}{20}{0}%
\regfield{UART\_RXD\_INV}{1}{19}{0}%
\regfield{UART\_TXFIFO\_RST}{1}{18}{0}%
\regfield{UART\_RXFIFO\_RST}{1}{17}{0}%
\regfield{UART\_IRDA\_EN}{1}{16}{0}%
\regfield{UART\_TX\_FLOW\_EN}{1}{15}{0}%
\regfield{UART\_LOOPBACK}{1}{14}{0}%
\regfield{UART\_IRDA\_RX\_INV}{1}{13}{0}%
\regfield{UART\_IRDA\_TX\_INV}{1}{12}{0}%
\regfield{UART\_IRDA\_WCTL}{1}{11}{0}%
\regfield{UART\_IRDA\_TX\_EN}{1}{10}{0}%
\regfield{UART\_IRDA\_DPLX}{1}{9}{0}%
\regfield{UART\_TXD\_BRK}{1}{8}{0}%
\regfield{UART\_SW\_DTR}{1}{7}{0}%
\regfield{UART\_SW\_RTS}{1}{6}{0}%
\regfield{UART\_STOP\_BIT\_NUM}{2}{4}{{1}}%
\regfield{UART\_BIT\_NUM}{2}{2}{{3}}%
\regfield{UART\_PARITY\_EN}{1}{1}{0}%
\regfield{UART\_PARITY}{1}{0}{0}%
\reglabel{Reset}\regnewline%
\begin{regdesc}\begin{reglist}
\item [UART\_TICK\_REF\_ALWAYS\_ON]This register is used to select the clock; 1: APB clock; 0: REF\_TICK. (R/W)
\item [UART\_DTR\_INV]Set this bit to invert the level of the UART DTR signal. (R/W)
\item [UART\_RTS\_INV]Set this bit to invert the level of the UART RTS signal. (R/W)
\item [UART\_TXD\_INV]Set this bit to invert the level of the UART TxD signal. (R/W)
\item [UART\_DSR\_INV]Set this bit to invert the level of the UART DSR signal. (R/W)
\item [UART\_CTS\_INV]Set this bit to invert the level of the UART CTS signal. (R/W)
\item [UART\_RXD\_INV]Set this bit to invert the level of the UART Rxd signal. (R/W)
\item [UART\_TXFIFO\_RST]Set this bit to reset the UART transmit-FIFO.
{\bf NOTICE}: UART2 doesn't have any register to reset Tx\_FIFO or Rx\_FIFO, and the UART1\_TXFIFO\_RST and UART1\_RXFIFO\_RST in UART1 may impact the functioning of UART2. Therefore, these two registers in UART1 should only be used when the Tx\_FIFO and Rx\_FIFO in UART2 do not have any data.
 (R/W)
\item [UART\_RXFIFO\_RST]Set this bit to reset the UART receive-FIFO. {\bf NOTICE}: UART2 doesn't have any register to reset Tx\_FIFO or Rx\_FIFO, and the UART1\_TXFIFO\_RST and UART1\_RXFIFO\_RST in UART1 may impact the functioning of UART2. Therefore, these two registers in UART1 should only be used when the Tx\_FIFO and Rx\_FIFO in UART2 do not have any data.
 (R/W)
\item [UART\_IRDA\_EN]Set this bit to enable the IrDA protocol. (R/W)
\item [UART\_TX\_FLOW\_EN]Set this bit to enable the flow control function for the transmitter. (R/W)
\item [UART\_LOOPBACK]Set this bit to enable the UART loopback test mode. (R/W)
\item [UART\_IRDA\_RX\_INV]Set this bit to invert the level of the IrDA receiver. (R/W)
\item [UART\_IRDA\_TX\_INV]Set this bit to invert the level of the IrDA transmitter. (R/W)
\item [UART\_IRDA\_WCTL] 1: The IrDA transmitter's 11th bit is the same as its 10th bit; 0: set IrDA transmitter's 11th bit to 0. (R/W)
\item [UART\_IRDA\_TX\_EN]This is the start enable bit of the IrDA transmitter. (R/W)
\item [UART\_IRDA\_DPLX]Set this bit to enable the IrDA loopback mode. (R/W)
\item [UART\_TXD\_BRK]Set this bit to enable the transmitter to send NULL, when the process of sending data is completed. (R/W)
\item [Continued on the next page...]
\end{reglist}\end{regdesc}
\end{register}
\addtocounter{Regfloat}{-1}
\begin{register}{H}{UART\_CONF0\_REG}{0x20}
\begin{regdesc}\begin{reglist}
\item [Continued from the previous page...]
\item [UART\_SW\_DTR]This register is used to configure the software DTR signal used in software flow control. (R/W)
\item [UART\_SW\_RTS] This bit is used in hardware flow control when UART\_RX\_FLOW\_EN is 0. Set this bit to drive the RTS (or rtsn\_out) signal low, and reset to drive the signal high. (R/W)
\item [UART\_STOP\_BIT\_NUM]This register is used to set the length of the stop bit.\\0: Invalid. No effect\\
1: 1 bit \\
2: 1.5 bits\\
3: 2 bits\\ (R/W)
\item [UART\_BIT\_NUM]This register is used to set the length of data; 0: 5 bits, 1: 6 bits, 2: 7 bits, 3: 8 bits. (R/W)
\item [UART\_PARITY\_EN]Set this bit to enable the UART parity check. (R/W)
\item [UART\_PARITY]This register is used to configure the parity check mode;  0: even, 1: odd. (R/W)
\end{reglist}\end{regdesc}
\end{register}

\begin{register}{H}{UART\_CONF1\_REG}{0x24}\label{regdesc:UARTCONF1REG}
\regfield{UART\_RX\_TOUT\_EN}{1}{31}{0}%
\regfield{UART\_RX\_TOUT\_THRHD}{7}{24}{0000000}%
\regfield{UART\_RX\_FLOW\_EN}{1}{23}{0}%
\regfield{UART\_RX\_FLOW\_THRHD}{7}{16}{{0x00}}%
\regfield{(reserved)}{1}{15}{0}%
\regfield{UART\_TXFIFO\_EMPTY\_THRHD}{7}{8}{{0x60}}%
\regfield{(reserved)}{1}{7}{0}%
\regfield{UART\_RXFIFO\_FULL\_THRHD}{7}{0}{{0x60}}%
\reglabel{Reset}\regnewline%
\begin{regdesc}\begin{reglist}
\item [UART\_RX\_TOUT\_EN]This is the enable bit for the UART receive-timeout function. (R/W)
\item [UART\_RX\_TOUT\_THRHD]\label{UARTRXTOUTTHRHD} This register is used to configure the UART receiver's timeout value when receiving a byte. When using APB\_CLK as the clock source, the register counts by UART baud cycle multiplied by 8. When using REF\_TICK as the clock source, the register counts by\\ UART baud cycle * 8 * (REF\_TICK frequency)/(APB\_CLK frequency). (R/W)
\item [UART\_RX\_FLOW\_EN]This is the flow enable bit of the UART receiver; 1: choose software flow control by configuring the sw\_rts signal; 0: disable software flow control. (R/W)
\item [UART\_RX\_FLOW\_THRHD]When UART\_RX\_FLOW\_EN is 1 and the receiver gets more data than its threshold value, the receiver produces an rtsn\_out signal that tells the transmitter to stop transferring data. The threshold value is (rx\_flow\_thrhd\_h3, rx\_flow\_thrhd). (R/W)
\item [UART\_TXFIFO\_EMPTY\_THRHD]\label{UARTTXFIFOEMPTYTHRHD}When the data amount in transmit-FIFO is less than its threshold value, it will produce a TXFIFO\_EMPTY\_INT\_RAW interrupt. The threshold value is (tx\_mem\_empty\_thrhd, txfifo\_empty\_thrhd). (R/W)
\item [UART\_RXFIFO\_FULL\_THRHD]\label{UARTRXFIFOFULLTHRHD} When the receiver gets more data than its threshold value, the receiver will produce an RXFIFO\_FULL\_INT\_RAW interrupt. The threshold value is (rx\_flow\_thrhd\_h3, rxfifo\_full\_thrhd). (R/W)
\end{reglist}\end{regdesc}
\end{register}


\begin{register}{H}{UART\_LOWPULSE\_REG}{0x28}\label{regdesc:UARTLOWPULSEREG}
\regfield{(reserved)}{12}{20}{000000000000}%
\regfield{UART\_LOWPULSE\_MIN\_CNT}{20}{0}{{0x0FFFFF}}%
\reglabel{Reset}\regnewline%
\begin{regdesc}\begin{reglist}
\item [UART\_LOWPULSE\_MIN\_CNT]This register stores the value of the minimum duration of the low-level pulse. It is used in the baud rate detection process. (RO)
\end{reglist}\end{regdesc}
\end{register}


\begin{register}{H}{UART\_HIGHPULSE\_REG}{0x2C}\label{regdesc:UARTHIGHPULSEREG}
\regfield{(reserved)}{12}{20}{000000000000}%
\regfield{UART\_HIGHPULSE\_MIN\_CNT}{20}{0}{{0x0FFFFF}}%
\reglabel{Reset}\regnewline%
\begin{regdesc}\begin{reglist}
\item [UART\_HIGHPULSE\_MIN\_CNT]This register stores  the value of the minimum duration of the high level pulse. It is used in baud rate detection process. (RO)
\end{reglist}\end{regdesc}
\end{register}


\begin{register}{H}{UART\_RXD\_CNT\_REG}{0x30}\label{regdesc:UARTRXDCNTREG}
\regfield{(reserved)}{22}{10}{0000000000000000000000}%
\regfield{UART\_RXD\_EDGE\_CNT}{10}{0}{{0x000}}%
\reglabel{Reset}\regnewline%
\begin{regdesc}\begin{reglist}
\item [UART\_RXD\_EDGE\_CNT]This register stores the count of the RxD edge change. It is used in the baud rate detection process. (RO)
\end{reglist}\end{regdesc}
\end{register}


\begin{register}{H}{UART\_FLOW\_CONF\_REG}{0x34}\label{regdesc:UARTFLOWCONFREG}
\regfield{(reserved)}{26}{6}{00000000000000000000000000}%
\regfield{UART\_SEND\_XOFF}{1}{5}{0}%
\regfield{UART\_SEND\_XON}{1}{4}{0}%
\regfield{UART\_FORCE\_XOFF}{1}{3}{0}%
\regfield{UART\_FORCE\_XON}{1}{2}{0}%
\regfield{UART\_XONOFF\_DEL}{1}{1}{0}%
\regfield{UART\_SW\_FLOW\_CON\_EN}{1}{0}{0}%
\reglabel{Reset}\regnewline%
\begin{regdesc}\begin{reglist}
\item [UART\_SEND\_XOFF]Hardware auto-clear; set to 1 to send Xoff char. (R/W)
\item [UART\_SEND\_XON]Hardware auto-clear; set to 1 to send Xon char. (R/W)
\item [UART\_FORCE\_XOFF]\label{regdesc:UARTFORCEXOFF} Set this bit to set the internal CTSn and stop the transmitter from sending data. (R/W)
\item [UART\_FORCE\_XON]\label{regdesc:UARTFORCEXON} Set this bit to clear the internal CTSn and enable the transmitter to continue sending data. (R/W)
\item [UART\_XONOFF\_DEL]Set this bit to remove the flow-control char from the received data. (R/W)
\item [UART\_SW\_FLOW\_CON\_EN]Set this bit to enable software flow control. It is used with register sw\_xon or sw\_xoff. (R/W)
\end{reglist}\end{regdesc}
\end{register}


\begin{register}{H}{UART\_SLEEP\_CONF\_REG}{0x38}\label{regdesc:UARTSLEEPCONFREG}
\regfield{(reserved)}{22}{10}{0000000000000000000000}%
\regfield{UART\_ACTIVE\_THRESHOLD}{10}{0}{{0x0F0}}%
\reglabel{Reset}\regnewline%
\begin{regdesc}\begin{reglist}
\item [UART\_ACTIVE\_THRESHOLD]\label{UARTACTIVETHRESHOLD}When the number of positive edges of RxD signal is larger than or equal to (UART\_ACTIVE\_THRESHOLD+2), the system emerges from Light-sleep mode and becomes active. (R/W)
\end{reglist}\end{regdesc}
\end{register}


\begin{register}{H}{UART\_SWFC\_CONF\_REG}{0x3C}\label{regdesc:UARTSWFCCONFREG}
\regfield{UART\_XOFF\_CHAR}{8}{24}{{0x013}}%
\regfield{UART\_XON\_CHAR}{8}{16}{{0x011}}%
\regfield{UART\_XOFF\_THRESHOLD}{8}{8}{{0x0E0}}%
\regfield{UART\_XON\_THRESHOLD}{8}{0}{{0x000}}%
\reglabel{Reset}\regnewline%
\begin{regdesc}\begin{reglist}
\item [UART\_XOFF\_CHAR]This register stores the Xoff flow control char. (R/W)
\item [UART\_XON\_CHAR]This register stores the Xon flow control char. (R/W)
\item [UART\_XOFF\_THRESHOLD]\label{UARTXOFFTHRESHOLD}When the data amount in receive-FIFO is more than what this register indicates, it will send an Xoff char, with uart\_sw\_flow\_con\_en set to 1. (R/W)
\item [UART\_XON\_THRESHOLD]\label{UARTXONTHRESHOLD}When the data amount in receive-FIFO is less than what this register indicates, it will send an Xon char, with uart\_sw\_flow\_con\_en set to 1. (R/W)
\end{reglist}\end{regdesc}
\end{register}


\begin{register}{H}{UART\_IDLE\_CONF\_REG}{0x40}\label{regdesc:UARTIDLECONFREG}
\regfield{(reserved)}{4}{28}{0000}%
\regfield{UART\_TX\_BRK\_NUM}{8}{20}{{0x00A}}%
\regfield{UART\_TX\_IDLE\_NUM}{10}{10}{{0x100}}%
\regfield{UART\_RX\_IDLE\_THRHD}{10}{0}{{0x100}}%
\reglabel{Reset}\regnewline%
\begin{regdesc}\begin{reglist}
\item [UART\_TX\_BRK\_NUM]This register is used to configure the number of zeros (0) sent, after the process of sending data is completed. It is active when txd\_brk is set to 1. (R/W)
\item [UART\_TX\_IDLE\_NUM]This register is used to configure the duration between transfers. (R/W)
\item [UART\_RX\_IDLE\_THRHD]When the receiver takes more time to receive Byte data than what this register indicates, it will produce a frame-end signal. (R/W)
\end{reglist}\end{regdesc}
\end{register}


\begin{register}{H}{UART\_RS485\_CONF\_REG}{0x44}\label{regdesc:UARTRS485CONFREG}
\regfield{(reserved)}{22}{10}{0000000000000000000000}%
\regfield{UART\_RS485\_TX\_DLY\_NUM}{4}{6}{0000}%
\regfield{UART\_RS485\_RX\_DLY\_NUM}{1}{5}{0}%
\regfield{UART\_RS485RXBY\_TX\_EN}{1}{4}{0}%
\regfield{UART\_RS485TX\_RX\_EN}{1}{3}{0}%
\regfield{UART\_DL1\_EN}{1}{2}{0}%
\regfield{UART\_DL0\_EN}{1}{1}{0}%
\regfield{UART\_RS485\_EN}{1}{0}{0}%
\reglabel{Reset}\regnewline%
\begin{regdesc}\begin{reglist}
\item [UART\_RS485\_TX\_DLY\_NUM]This register is used to delay the transmitter's internal data signal. (R/W)
\item [UART\_RS485\_RX\_DLY\_NUM]This register is used to delay the receiver's internal data signal. (R/W)
\item [UART\_RS485RXBY\_TX\_EN]1: enable the RS-485 transmitter to send data, when the RS-485 receiver line is busy; 0: the RS-485 transmitter should not send data, when its receiver is busy. (R/W)
\item [UART\_RS485TX\_RX\_EN]Set this bit to enable the transmitter's output signal loop back to the receiver's input signal.  (R/W)
\item [UART\_DL1\_EN]Set this bit to delay the STOP bit by 1 bit. (R/W)
\item [UART\_DL0\_EN]\label{UARTDL0EN}Set this bit to delay the STOP bit by 1 bit after DL1. (R/W)
\item [UART\_RS485\_EN]Set this bit to choose the RS-485 mode. (R/W)
\end{reglist}\end{regdesc}
\end{register}


\begin{register}{H}{UART\_AT\_CMD\_PRECNT\_REG}{0x48}\label{regdesc:UARTATCMDPRECNTREG}
\regfield{(reserved)}{8}{24}{00000000}%
\regfield{UART\_PRE\_IDLE\_NUM}{24}{0}{{0x0186A00}}%
\reglabel{Reset}\regnewline%
\begin{regdesc}\begin{reglist}
\item [UART\_PRE\_IDLE\_NUM]This register is used to configure the idle-time duration before the first at\_cmd is received by the receiver. When the duration is less than what this register indicates, it will not take the next data received as an at\_cmd char. (R/W)
\end{reglist}\end{regdesc}
\end{register}


\begin{register}{H}{UART\_AT\_CMD\_POSTCNT\_REG}{0x4c}\label{regdesc:UARTATCMDPOSTCNTREG}
\regfield{(reserved)}{8}{24}{00000000}%
\regfield{UART\_POST\_IDLE\_NUM}{24}{0}{{0x0186A00}}%
\reglabel{Reset}\regnewline%
\begin{regdesc}\begin{reglist}
\item [UART\_POST\_IDLE\_NUM]This register is used to configure the duration between the last at\_cmd and the next data. When the duration is less than what this register indicates, it will not take the previous data as an at\_cmd char. (R/W)
\end{reglist}\end{regdesc}
\end{register}


\begin{register}{H}{UART\_AT\_CMD\_GAPTOUT\_REG}{0x50}\label{regdesc:UARTATCMDGAPTOUTREG}
\regfield{(reserved)}{8}{24}{00000000}%
\regfield{UART\_RX\_GAP\_TOUT}{24}{0}{{0x0001E00}}%
\reglabel{Reset}\regnewline%
\begin{regdesc}\begin{reglist}
\item [UART\_RX\_GAP\_TOUT]\label{UARTRXGAPTOUT} This register is used to configure the interval between the at\_cmd chars. When the interval is greater than the value of this register, it will not take the data as continuous at\_cmd chars. The register should be configured to more than half of the baud rate. (R/W)
\end{reglist}\end{regdesc}
\end{register}


\begin{register}{H}{UART\_AT\_CMD\_CHAR\_REG}{0x54}\label{regdesc:UARTATCMDCHARREG}
\regfield{(reserved)}{16}{16}{0000000000000000}%
\regfield{UART\_CHAR\_NUM}{8}{8}{{0x003}}%
\regfield{UART\_AT\_CMD\_CHAR}{8}{0}{{0x02B}}%
\reglabel{Reset}\regnewline%
\begin{regdesc}\begin{reglist}
\item [UART\_CHAR\_NUM]This register is used to configure the number of continuous at\_cmd chars received by the receiver. (R/W)
\item [UART\_AT\_CMD\_CHAR]This register is used to configure the content of an at\_cmd char. (R/W)
\end{reglist}\end{regdesc}
\end{register}


\begin{register}{H}{UART\_MEM\_CONF\_REG}{0x58}\label{regdesc:UARTMEMCONFREG}
\regfield{(reserved)}{1}{31}{0}%
\regfield{UART\_TX\_MEM\_EMPTY\_THRHD}{3}{28}{{0x0}}%
\regfield{UART\_RX\_MEM\_FULL\_THRHD}{3}{25}{{0x0}}%
\regfield{UART\_XOFF\_THRESHOLD\_H2}{2}{23}{{0x0}}%
\regfield{UART\_XON\_THRESHOLD\_H2}{2}{21}{{0x0}}%
\regfield{UART\_RX\_TOUT\_THRHD\_H3}{3}{18}{{0x0}}%
\regfield{UART\_RX\_FLOW\_THRHD\_H3}{3}{15}{{0x0}}%
\regfield{(reserved)}{4}{11}{0000}%
\regfield{UART\_TX\_SIZE}{4}{7}{{0x01}}%
\regfield{UART\_RX\_SIZE}{4}{3}{{0x01}}%
\regfield{(reserved)}{2}{1}{00}%
\regfield{UART\_MEM\_PD}{1}{0}{0}%
\reglabel{Reset}\regnewline%
\begin{regdesc}\begin{reglist}
\item [UART\_TX\_MEM\_EMPTY\_THRHD]Refer to the description of \hyperref[UARTTXFIFOEMPTYTHRHD]{TXFIFO\_EMPTY\_THRHD}. (R/W)
\item [UART\_RX\_MEM\_FULL\_THRHD]Refer to the description of \hyperref[UARTRXFIFOFULLTHRHD]{RXFIFO\_FULL\_THRHD}. (R/W)
\item [UART\_XOFF\_THRESHOLD\_H2]Refer to the description of \hyperref[UARTXOFFTHRESHOLD]{UART\_XOFF\_THRESHOLD}. (R/W)
\item [UART\_XON\_THRESHOLD\_H2]Refer to the description of \hyperref[UARTXONTHRESHOLD]{UART\_XON\_THRESHOLD}. (R/W)
\item [UART\_RX\_TOUT\_THRHD\_H3]Refer to the description of \hyperref[UARTRXTOUTTHRHD]{RX\_TOUT\_THRHD}. (R/W)
\item [UART\_RX\_FLOW\_THRHD\_H3]Refer to the description of \hyperref[UARTRXFIFOFULLTHRHD]{RX\_FLOW\_THRHD}. (R/W)
\item [UART\_TX\_SIZE]This register is used to configure the amount of memory allocated to the transmit-FIFO. The default number is 128 bytes. (R/W)
\item [UART\_RX\_SIZE]This register is used to configure the amount of memory allocated to the receive-FIFO. The default number is 128 bytes. (R/W)
\item [UART\_MEM\_PD]Set this bit to power down the memory. When the reg\_mem\_pd register is set to 1 for all UART controllers, Memory will enter the low-power mode. (R/W)
\end{reglist}\end{regdesc}
\end{register}

\begin{register}{H}{UART\_MEM\_TX\_STATUS\_REG}{0x5c}\label{regdesc:UARTMEMTXSTATUSREG}
\regfield{(reserved)}{8}{24}{0}%
\regfield{UART\_MEM\_TX\_WR\_ADDR}{11}{13}{{0}}%
\regfield{UART\_MEM\_TX\_RD\_ADDR}{11}{2}{0}%
\regfield{(reserved)}{2}{0}{{0}}%
\reglabel{Reset}\regnewline%
\begin{regdesc}\begin{reglist}
\item [UART\_MEM\_TX\_WR\_ADDR] Represents the offset address to write TX FIFO. (RO)
\item [UART\_MEM\_TX\_RD\_ADDR] Represents the offset address to read TX FIFO. (RO)
\end{reglist}\end{regdesc}
\end{register}


\begin{register}{H}{UART\_MEM\_RX\_STATUS\_REG}{0x0060}\label{regdesc:UARTMEMRXSTATUSREG}
\regfield{(reserved)}{8}{24}{0}%
\regfield{UART\_MEM\_RX\_WR\_ADDR}{11}{13}{{0}}%
\regfield{UART\_MEM\_RX\_RD\_ADDR}{11}{2}{0}%
\regfield{(reserved)}{2}{0}{{0}}%
\reglabel{Reset}\regnewline%
\begin{regdesc}\begin{reglist}
\item [UART\_MEM\_RX\_RD\_ADDR] Represents the offset address to read RX FIFO. (RO)
\item [UART\_MEM\_RX\_WR\_ADDR] Represents the offset address to write RX FIFO. (RO)
\end{reglist}\end{regdesc}
\end{register}

\begin{register}{H}{UART\_MEM\_CNT\_STATUS\_REG}{0x64}\label{regdesc:UARTMEMCNTSTATUSREG}
\regfield{(reserved)}{26}{6}{00000000000000000000000000}%
\regfield{UART\_TX\_MEM\_CNT}{3}{3}{000}%
\regfield{UART\_RX\_MEM\_CNT}{3}{0}{000}%
\reglabel{Reset}\regnewline%
\begin{regdesc}\begin{reglist}
\item [UART\_TX\_MEM\_CNT]Refer to the description of \hyperref[UARTTXFIFOCNT]{TXFIFO\_CNT}. (RO)
\item [UART\_RX\_MEM\_CNT]Refer to the description of \hyperref[UARTRXFIFOCNT]{RXFIFO\_CNT}. (RO)
\end{reglist}\end{regdesc}
\end{register}


\begin{register}{H}{UART\_POSPULSE\_REG}{0x68}\label{regdesc:UARTPOSPULSEREG}
\regfield{(reserved)}{12}{20}{000000000000}%
\regfield{UART\_POSEDGE\_MIN\_CNT}{20}{0}{{0x0FFFFF}}%
\reglabel{Reset}\regnewline%
\begin{regdesc}\begin{reglist}
\item [UART\_POSEDGE\_MIN\_CNT]This register stores the count of RxD positive edges. It is used in the autobaud detection process. (RO)
\end{reglist}\end{regdesc}
\end{register}


\begin{register}{H}{UART\_NEGPULSE\_REG}{0x6c}\label{regdesc:UARTNEGPULSEREG}
\regfield{(reserved)}{12}{20}{000000000000}%
\regfield{UART\_NEGEDGE\_MIN\_CNT}{20}{0}{{0x0FFFFF}}%
\reglabel{Reset}\regnewline%
\begin{regdesc}\begin{reglist}
\item [UART\_NEGEDGE\_MIN\_CNT]This register stores the count of RxD negative edges. It is used in the autobaud detection process. (RO)
\end{reglist}\end{regdesc}
\end{register}

